\section{Gap Synthesis and Conceptual Framework}\label{sec:gap-synthesis}

The review in this chapter converges on three findings. First, the
language-model copilots surveyed \parencite{deng2024pentestgpt,happe2023pwned}
generate plausible security actions but do not validate them for machine
consumption: their output is produced for a human reviewer, and correctness is
handled informally (Section~\ref{sec:related-research}). Second, applications
that integrate a model with tooling are exploitable through adversarial and
indirect-injection utterances, the surveyed defences are individually
insufficient \parencite{liu2023promptinjection,greshake2023indirect}, and the
existing-systems survey shows that current guardrail techniques act as
post-hoc application-layer checks rather than as a deterministic validation
stage between generation and dispatch (Section~\ref{sec:existing-systems}).
Third, fabrication-detection and agreement methods
\parencite{manakul2023selfcheckgpt,kuhn2023semantic,farquhar2024semantic,wang2023selfconsistency}
are available in principle, but they assume sampling or interactive generation
whose cost collides with the offline, CPU-only, single-inference budget adopted
here (Section~\ref{sec:domain-foundations}).

The gap that follows from these three findings is the combination that the
existing-systems survey leaves open: no surveyed system delivers a dedicated,
isolated natural-language understanding layer that converts free-form
offensive-security instructions into validated, schema-constrained, authorised
machine contracts while running offline, retrieval-free, on a single inference
budget. This dissertation addresses exactly that combination by designing and
\emph{measuring} a guarded intent-recognition engine.

\subsection*{Conceptual framework}

The design chapters build on three principles established in
Section~\ref{sec:domain-foundations}.

\emph{First}, the language model is treated as a \emph{stochastic proposition
generator}: a single generated response is a sample weighted by the model's
beliefs rather than a deterministic verdict
(Section~\ref{sec:domain-foundations}). Nothing the model emits is trusted until
it satisfies checkable predicates; sampling-based agreement
\parencite{wang2023selfconsistency} provides a confidence signal, but the
generation itself is always a candidate rather than a decision.

\emph{Second}, responsibility is divided along the decidability boundary
identified in Section~\ref{sec:domain-foundations}: syntax, format and
authorisation are decidable predicates (schema, host and CIDR syntax, port
ranges, network and scope policy, role permission), while only meaning is
delegated to the model. The guarded pipeline therefore couples a small
instruction-tuned model with deterministic stages---schema validation, regex and
network validation, scope enforcement, role-based access control---whose purpose
is to contain model error before it can propagate to a planner or executor
(Chapter~\ref{ch:product-design}).

\emph{Third}, the design enforces a \emph{safety invariant}: an intent contract
is accepted only when it satisfies the conjunction of the deterministic security
predicates---schema-valid, network-authorized, in-scope, role-permitted and
non-adversarial. The model may propose any candidate, but it cannot force
acceptance of a candidate that violates any predicate in the conjunction. This
converts the open problem of making a single language model trustworthy into the
closed problem of classifying a single utterance under a known policy---the
reduction that justifies acting on model-produced security instructions in an
automated toolchain.

On top of these principles the framework adopts the bounded-sampling resolution
of the cost tension from Section~\ref{sec:domain-foundations}: a small, bounded
number of responses, a cheap embedding-based agreement gate, and a deterministic
escalation path (the hallucination validator) that does not depend on a large
verifier model (Section~\ref{sec:related-research}).

\subsection*{From framework to measurement}

The framework yields three measurable claims, each mapped to the corresponding
research question (Section~\ref{sec:research-questions}):
\begin{itemize}
\item \textbf{Pipeline value (RQ1).} The guarded pipeline raises the
probability that model output satisfies the full conjunction of security
predicates relative to raw generation; this is quantified by comparing intent
accuracy, flaw-catch rate and false-positive rate with and without the guards.
\item \textbf{Layer contribution (RQ2).} The contribution of each guard layer is
quantified by ablating layers in turn against a constant baseline, isolating the
largest safety contributor and the layers that trade safety for usability.
\item \textbf{Generality (RQ3).} The identical guardrail stack is applied to two
distinct offline backends to determine how much of the measured safety is
backend-independent rather than an artefact of a single model.
\end{itemize}
The evaluation protocol that operationalizes these questions is developed in
Chapter~\ref{ch:methodology}, the implementation and automated test suite in
Chapter~\ref{ch:implementation-testing}, and the measurements themselves in
Chapter~\ref{ch:results}. Together with the preceding sections, this concludes
the review, sets out the conceptual framework that guides the design of the
Intent Recognition Engine in Chapter~\ref{ch:product-design}.

\TODO{At the final AUDIT step, cross-check the informal safety invariant stated
here against the full formal definition to be developed in
Chapter~\ref{ch:product-design}.}